\section{Details of Causal and Context Interventions}
\label{app:ablation}
The native ablation evaluation uses 253 frontier assignments. Both intervention types begin at the final prompt token and continue at subsequent query positions during decoding; earlier query rows remain unmodified. For span restriction, the blocked key interval begins one token before the frontier description and extends through the prompt's final assignment token. Heads in layers 20--63 outside the spared set are blocked from that interval. For sink restriction, selected heads attend only to key position zero. The latter produces the sink value contribution and is not mathematically identical to zeroing a head.

The 122-head core and 250-head set differ in both membership and depth coverage. Random controls sample from the corresponding depth pool, not a matched count at each individual layer. For a set containing heads below layer 20, those heads are outside the span-restriction band. Thus total selected count and number of spared heads inside that band must be distinguished. The experiments test these specified interventions, not arbitrary removal of the entire complement or isolation of the selected heads from the rest of the network.

\subsection{Task and query transfer}
\label{app:task_transfer}
Task-span KV transfer uses per-query-head output routing to avoid transferring to every query head sharing a KV head. Task spans are padded to equal length and source keys are re-rotated to destination positions; equal lengths do not imply identical absolute starts. The KV patch is active across prefill positions and decoding. Query transfer patches the specified pivot tokens and reuses the source pivot query during the steering window (32 forward passes in the evaluation CLI, including prefill). The reported transfer configurations also amplify selected OV contributions with coefficient 2; this differs from the GoalScope probe's coefficient 0.1.

The principal multi-predicate 250-head runs report approximately 94\% source-directed generation for query transfer ($n=2048$) and 99\% for KV transfer at the assignment pivot ($n=512$), spanning four datatypes and four language directions. The latter uses no non-goal-head ablation. These rates count source-or-both oracle labels and do not imply perfect program correctness. GoalScope's patched source readout is a separate outcome. Q-filter examples exhibit readout/behavior disagreement; aggregate readout success should not be treated as a substitute for implementation success.

The qualitative task-order example in the main text follows the saved figure with unequal source/destination progress. It should not be described as an identical-progress experiment. Broader conclusions about query content require a systematic crossed design. We make no claim that these interventions demonstrate a particular set-subtraction operation.

\section{OVReader}
\label{app:ovreader}
For a residual state after source layer $L$, OVReader applies each selected head's value map and output map, including its pre-attention normalization gain. It unit-normalizes each transformed vector and averages them. This operator is based on model weights, but is not a single linear projection because of the per-head normalization. The aggregate steers a separate GoalScope probe with norm-relative coefficient .15; the selected-head-output comparison uses .1.

The evaluation samples 64 examples per datatype, with five capture positions and layer sweeps through the final layer. At the assignment operator, readout can approach the head-output reference around L40; pooled curves generally peak near L48--52. A decline at the final layer indicates reduced accessibility through this particular transformation and probe, not absence of all subgoal information. Implementation-transfer scores were rejudged with a trivial-output sanitizer; existing records should not be described as a freshly resampled corpus excluding all trivial target predicates.
